\begin{abstract}
Visual perturbations that leave a task unchanged can sharply reduce the success of vision-language-action (VLA) policies. Making these policies robust requires knowing what a perturbation changes inside them and how to judge a correction. How much a perturbation changes features, commands, or motion (its \emph{sensitivity}) is the natural measure, yet it need not tell how much success the perturbation costs (its \emph{fragility}). We trace perturbations through frozen $\pi_{0.5}$ and $\pi_0$ on LIBERO and LIBERO-Plus, from features to commands, motion, and task outcomes. Inside $\pi_{0.5}$, different perturbations change the features along 16 shared directions per layer: on held-out tasks and perturbation families, these capture 36.5--83.6\% of the variation in the changes, against at most 10.2\% for other sets of 16 directions. They give one common place to intervene. Copying clean values back along them, using a clean observation of the same state, reduces the error of the first five commands by 24\% under camera changes and 5\% under sensor noise, and task success is higher as well. These policies generate actions from random noise, so two unperturbed runs already move differently, which gives a yardstick: the clean rerun. Perturbation families that move the robot more tend to cost more success, but within a family the size of the early motion change barely tells which runs fail. Under sensor noise, $\pi_0$ deviates less from its clean run than a clean rerun does, and a rerun never changes the outcome, yet the perturbation costs 14.3 points of success. Net success also hides lost and gained successes. Robustness, and any internal correction, should therefore be judged by paired task outcomes against clean reruns.
\end{abstract}
